\subsection{Full subalgebras}

Let A be a unital algebra over K.

DEFINITION 3. --- A full subalgebra of A is a unital subalgebra B such that every element of B which is invertible in A is invertible in B.

In other words, B is a full subalgebra of A if and only if \(\mathrm{Sp}_{\mathrm{B}}(x)=\mathrm{Sp}_{\mathrm{A}}(x)\) for all \(x\in B\) .

The intersection of a family of full subalgebras of A is a full subalgebra of A.

Let M be a subset of A. The intersection of the full subalgebras of A containing M is the smallest full subalgebra of A containing M.

It is called the full subalgebra of A generated by M. The commutant M′ of M in A is a full subalgebra of A (for, if x is invertible in A and commutes with M, then \(x^{-1}\) commutes with M). Therefore, the bicommutant M'' of M is a full subalgebra of A that contains the full subalgebra of A generated by M.

If the elements of M are pairwise permutable, we have \(M \subset M'\) and \(M'' \subset M'''\) ; the algebra \(M''\) is therefore commutative, and the same is then true of the full subalgebra generated by M.

A maximal commutative subalgebra of A is a full subalgebra, because it is equal to its commutant.

Lemma 2. --- Let \((x_{\lambda})_{\lambda \in \Lambda}\) be a family of pairwise commuting elements of A. The full subalgebra generated by \((x_{\lambda})\) is the set of elements of the form \(\mathrm{R}((x_{\lambda}))\) , where \(\mathrm{R} \in \mathrm{K}((\mathrm{X}_{\lambda}))\) ranges over the set of rational fractions in which the family \((x_{\lambda})\) is substitutable.

Let B be the full subalgebra of A generated by the family \((x_{\lambda})\) , and denote by \(B_{1}\) the set of elements of the form \(\mathrm{R}((x_{\lambda}))\) , where \(\mathrm{R}((\mathrm{X}_{\lambda}))\) is a rational fraction in which \((x_{\lambda})\) is substitutable. Explicitly, \(B_{1}\) is the set of elements of A of the form \(\mathrm{P}((x_{\lambda}))\mathrm{Q}((x_{\lambda}))^{-1}\) , where \(\mathrm{P},\mathrm{Q}\in\mathrm{K}[(\mathrm{X}_{\lambda})]\) and \(\mathrm{Q}((x_{\lambda}))\) is invertible in A. The set \(B_{1}\) is a unital subalgebra of A containing the family \((x_{\lambda})\) . It is a full subalgebra: if \(\mathrm{P}((x_{\lambda}))\mathrm{Q}((x_{\lambda}))^{-1}\) is invertible in A, then \(\mathrm{P}((x_{\lambda}))\) is invertible in A and the inverse \(\mathrm{Q}((x_{\lambda}))\mathrm{P}((x_{\lambda}))^{-1}\) of \(\mathrm{P}((x_{\lambda}))\mathrm{Q}((x_{\lambda}))^{-1}\) belongs to \(B_{1}\) . We therefore have \(B\subset B_{1}\) . On the other hand, if \(P,Q\in K[(X_{\lambda})]\) , and if \(\mathrm{Q}((x_{\lambda}))\) is invertible in A, then \(\mathrm{P}((x_{\lambda}))\in\mathrm{B}\) and \(\mathrm{Q}((x_{\lambda}))\in\mathrm{B}\) , hence \(\mathrm{Q}((x_{\lambda}))^{-1}\in\mathrm{B}\) and \(\mathrm{P}((x_{\lambda}))\mathrm{Q}((x_{\lambda}))^{-1}\in\mathrm{B}\) , hence \(B_{1}\subset B\) .

